\chapter{Conclusion}
\label{ch:conclusion}

This thesis revisited class-incremental learning (CIL) in the context of modern visual foundation models, arguing that once representations are frozen and sufficiently discriminative, the central challenge of CIL shifts from protecting network parameters against interference to properly modelling the statistical structure of a stationary embedding space. We formalised this shift by formulating incremental learning as the recursive, gradient-free estimation of class-conditional Gaussian densities in the embedding space of a frozen DINOv2 backbone. Each class is represented entirely by its sufficient statistics---a mean vector and a covariance matrix---which are updated in closed form as new samples arrive, without backpropagation, replay buffers, exemplar storage, or prompt parameters, and without ever modifying the statistics of a previously learned class. This design eliminates catastrophic forgetting by construction rather than by mitigation, and it is the structural reason why the forgetting scores reported throughout Chapter~\ref{ch:results} reflect increasing classification difficulty as the class set grows rather than any degradation of previously acquired knowledge.

Empirically, this conceptually simple formulation is competitive with the strongest prompt-based state of the art. Our full-covariance and mixture-of-Gaussian configurations match ConvPrompt on Split CIFAR-100 ($88.87\,\%$) without any trainable parameters or optimisation during incremental learning, and the Full-Cov and Hybrid-Cov estimators outperform every prompt-based baseline reported on Split Caltech-256 ($94.93\,\%$ versus the best prompt-based result of $92.50\,\%$). This level of accuracy is achieved at a computational cost for the classification stage below $0.3$ GFLOPs in the worst case---more than fifty times smaller than the $\approx 17.56$ GFLOPs required for a single forward pass of the DINOv2 backbone---so that the probabilistic classification layer introduces only marginal overhead relative to feature extraction, and requires none of the gradient updates that prompt-based competitors perform at every incremental task.

Beyond the accuracy figures, this thesis provided a rigorous statistical validation of the assumption that underlies the entire framework. Combining moment-based diagnostics with formal hypothesis tests---the D'Agostino--Pearson omnibus test for univariate normality and a diagonal adaptation of the Henze--Zirkler test for multivariate normality under severe rank deficiency---confirmed that the marginal feature distributions of a self-supervised foundation encoder are consistently, if not perfectly, Gaussian across three visually heterogeneous benchmarks: $70$--$81\,\%$ of all (class, dimension) pairs pass the strict D'Agostino--Pearson test, and the diagonal multivariate Henze--Zirkler statistic converges tightly to its Gaussian reference value of $1$ on every dataset. The complementary Bhattacharyya-distance analysis established a direct, quantitative link between inter-class separability and the achievable classification accuracy: over $99.8\,\%$ of all class pairs on every benchmark exceed a separability threshold of $D_B=50$, while the small minority of pairs that fall below $D_B=20$---semantic or visual near-synonyms such as \textit{boy}/\textit{girl}, \textit{bow\_tie}/\textit{scarf}, and \textit{top-hat}/\textit{cowboy-hat}---account for essentially the entirety of the residual classification error identified in the confusion-pattern analysis of Chapter~\ref{ch:analysis}. This is, to our knowledge, one of the first attempts in the class-incremental learning literature to connect a formal probabilistic error bound directly to the specific, named class pairs that a deployed classifier will actually confuse.

This thesis further showed that the reliability of full-rank covariance modelling is itself a matter of numerical care, not only of statistical assumption. Computing the precision matrix and log-determinant of the Full-Cov and Hybrid-Cov estimators via a Cholesky-based decomposition, rather than via an eigendecomposition with a hard eigenvalue floor, avoids a substantial artefact that would otherwise depress the Full-Cov accuracy on Split Caltech-256 by nearly five accuracy points and produce misleadingly optimistic results on the 30-task Joint benchmark. The results present a clear and interpretable picture: full-rank covariance estimation is highly beneficial whenever the per-class sample budget is adequate and homogeneous across classes, but the Joint benchmark's collapse of Full-Cov and Hybrid-Cov to $59.35\,\%$ is not a weakness of the estimator---it is a confirmation that covariance parameterisation must be matched to the sample-size regime of the deployment scenario, precisely as the Shared-covariance variants, which pool statistics across classes of heterogeneous rank, demonstrate by becoming the strongest estimators in that same setting ($81.49\,\%$).

These results, taken together, support the view that in the visual foundation-model era, the central challenge of class-incremental learning is no longer protecting network parameters from catastrophic interference but correctly modelling and exploiting the statistical geometry of a stable representation space. The success of a classical, decades-old statistical tool---the multivariate Gaussian classifier---in this modern setting is not a coincidence: it is made possible by the qualitative change that large-scale self-supervised pretraining has produced in the topology of visual embedding spaces. By decoupling representational plasticity from semantic categorisation, stable and well-behaved foundation features eliminate representation drift by construction, allowing lifelong visual recognition to be reframed as an ongoing, closed-form problem of probability density estimation rather than an ongoing problem of parameter optimisation.

\paragraph{Limitations.} A key limitation of the present work follows directly from the discussion of Chapter~\ref{ch:analysis} and deserves explicit restatement here: the framework's residual error is structurally bounded by the semantic resolution of the upstream encoder. No choice of covariance estimator or mixture model evaluated in this thesis can fully resolve a class pair whose Bhattacharyya distance and Euclidean mean separation are both near zero, since such a pair reflects a genuine ambiguity in the frozen feature space rather than a deficiency of the downstream classifier.

\section{Future Work}

Several natural extensions of the framework deserve further investigation, building on the longer-term research directions identified in this thesis.

\textbf{Targeted Representation Alignment for Lingering Confusion Sinks.} 
Our failure diagnostics isolated a minute, semantically coherent subset of class pairs (e.g., \textit{boy}/\textit{girl} or \textit{top-hat}/\textit{cowboy-hat}) that account for almost the entire residual error of the system due to their near-zero mean separation ($d_E \to 0$). Because this crowding is an inherent property of the upstream embedding space, it cannot be resolved by downstream parametric tuning. Future research should investigate lightweight, targeted representation alignment. Instead of re-training the entire deep backbone, a contrastive or metric-learning objective can be applied exclusively to the low-$D_B$ class pairs flagged by our statistical diagnostics. This targeted adjustment can local-stretch the decision surface for highly ambiguous categories, maximizing class separability without disrupting the global stationarity of the remaining representation space.

\textbf{Structured Covariance Parameterisation Under Highly Heterogeneous Sample Budgets.} 
The Hybrid-Cov estimator successfully decouples the Mahalanobis shape component from the log-normalisation volume term, providing a stable baseline on homogeneous streams. However, as demonstrated by the Joint benchmark collapse, unconstrained second-order parameters break down when forced to compare data-rich classes directly against sample-starved categories. To alleviate this sensitivity, future work should develop structured precision matrices that dynamically adapt their complexity to individual sample counts.

\textbf{Probabilistic Uncertainty Calibration.} 
Our calibration metrics reveal that full-covariance parameterisations suffer from structural over-confidence, as the estimated distribution volumes tend to collapse around scarce sample clusters. This miscalibration reduces the trustworthiness of downstream probability scores in safety-critical deployment environments. Future iterations should implement post-hoc calibration layers designed for closed-form streams, such as multi-class temperature scaling or conformal prediction frameworks. By computing distribution-free prediction intervals based on historical validation residuals, the system can provide mathematically guaranteed error bounds alongside its standard maximum-a-posteriori assignments.


\textbf{Gradient-Free Spatial Object Detection.} 
A natural functional scaling for this architecture is extending its reach from global scene classification to localized object detection. Traditional multi-stage object detectors rely on intensive backpropagation to optimize region proposal networks (RPNs) and bounding-box regressors, creating severe computational bottlenecks. Future work should bypass these requirements by exploiting the localized patch tokens natively generated during a single forward pass of the DINOv2 backbone. By applying simple mathematical decompositions—such as localized principal component analysis or spatial attention routing—directly onto these frozen spatial features, areas of interest can be isolated and classified in a single, gradient-free pass, preserving the lightweight, deterministic philosophy of this thesis.

\textbf{Autonomous Open-World Discovery.} 
Because this framework computes closed-form sufficient statistics natively, it provides a well-behaved predictive entropy and Log-Gap score for every incoming sample. This capability can be leveraged to shift the system from a closed-set classifier toward an open-world discovery framework. Since the uncertainty boundaries of the Diag-Cov branch remain structurally stable when new categories are introduced, establishing an adaptive out-of-distribution (OOD) threshold would allow the model to autonomously isolate unknown entities. Once flagged, these unfamiliar clusters can be accumulated and parameterized as new Gaussian classes in closed form, expanding the label space without requiring external supervision or manual task boundaries.

